\documentclass[12pt,a4paper]{article}
\usepackage[utf8]{inputenc}
\usepackage[T1]{fontenc}
\usepackage{geometry}
\geometry{margin=0.9in}
\usepackage{booktabs}
\usepackage{tabularx}
\usepackage{hyperref}
\usepackage{xcolor}
\usepackage{enumitem}

\definecolor{primary}{HTML}{1B365D}
\definecolor{accent}{HTML}{2B6CB0}

\hypersetup{
    colorlinks=true,
    linkcolor=primary,
    urlcolor=accent,
    pdftitle={Section 02: Literature Survey & State-of-the-Art Review},
    pdfauthor={Anurag Saha, Ayan Karak}
}

\begin{document}

\begin{center}
    {\Large \textbf{\color{primary} JALPAIGURI GOVERNMENT ENGINEERING COLLEGE (JGEC)}}\\[0.3em]
    {\normalsize Department of Information Technology $\bullet$ 5th Semester Official Project Documentation}\\[1em]
    \rule{\linewidth}{1.5pt}\\[0.5em]
    {\LARGE \textbf{AUTOMATED SKIN DISEASE CLASSIFICATION USING DEEP LEARNING}}\\[0.3em]
    {\large \textit{Section 02: Literature Survey \& State-of-the-Art Review}}\\[0.3em]
    \rule{\linewidth}{1.5pt}\\[0.8em]
    \textbf{Submitted by:} Anurag Saha (Roll No: \texttt{24101106037}) \& Ayan Karak (Roll No: \texttt{24101106013})\\[0.3em]
    \textbf{Submitted to \& Supervised by:} Dr. Somnath Dhibar\\[0.3em]
    \textbf{Source Code Repository:} \url{https://github.com/sahaanurag207-oss/SKIN-DISEASE-CLASSIFICATION-FROM-IMAGES}
\end{center}

\vspace{0.8em}

\section{Literature Survey \& State-of-the-Art Review}
Automated dermatoscopic image analysis has evolved rapidly over the past decade, progressing from hand-crafted ABCD rule feature extractors to deep convolutional neural networks (CNNs), hierarchical Vision Transformers (ViTs), and modernized pure-convolutional architectures. A systematic review of foundational and contemporary literature establishes the methodological motivation for our design choices developed under the mentorship of \textbf{Dr. Somnath Dhibar}.

\subsection{Evolution of Deep Learning in Dermatoscopy}
\begin{itemize}[leftmargin=*]
    \item \textbf{Landmark CNN Feasibility (Esteva et al., Nature, 2017):} Demonstrated dermatologist-level skin cancer classification using an Inception-v3 architecture trained on 129,450 clinical and dermatoscopic images. While proving deep neural networks can match board-certified specialists, the reliance on proprietary data highlighted the urgent need for standardized public benchmarks.
    \item \textbf{Standardization via HAM10000 (Tschandl et al., Scientific Data, 2018):} Released 10,015 multi-source dermatoscopic images across seven pigmented lesion categories with over 50\% histopathologic confirmation. Crucially, Tschandl et al. documented that multiple images of the exact same physical lesion (captured at varying magnifications or angles) exist in the dataset and must be tracked via the \texttt{lesion\_id} metadata field.
    \item \textbf{ISIC Grand Challenge Ensembles (Codella et al., 2018; Gessert et al., 2020):} Top-performing entries in the ISIC 2018 and 2019 challenges established that transfer learning from large-scale natural image datasets, heavy geometric/color augmentation, label smoothing, and Test-Time Augmentation (TTA) consistently outperform single-crop inference on imbalanced dermatoscopy data.
    \item \textbf{Modernized Pure ConvNets --- ConvNeXt (Liu et al., CVPR, 2022):} Re-examined standard ResNet design spaces through the lens of Swin Transformers, introducing $7\times 7$ depthwise convolutions, inverted bottlenecks, LayerNorm, GELU activations, and a 1:1:3:1 stage compute ratio. ConvNeXt matches or exceeds Vision Transformer accuracy while preserving convolutional inductive biases crucial for medical datasets of moderate scale ($\sim$10k images).
    \item \textbf{Lesion-Level Data Leakage in HAM10000 Literature:} Reproducibility audits revealed that naive random or image-level stratified splits artificially inflate HAM10000 accuracy above 93--96\% because duplicate views of the same lesion appear in both training and test sets. Enforcing \texttt{lesion\_id} group partitioning (\texttt{StratifiedGroupKFold}) eliminates this leakage, yielding realistic, clinically trustworthy generalization benchmarks.
\end{itemize}

\subsection{Comparative Synthesis of Prior Approaches vs. Proposed Framework}
\begin{table}[h!]
\centering
\renewcommand{\arraystretch}{1.25}
\begin{tabularx}{\textwidth}{l l l X}
\toprule
\textbf{Study / Paradigm} & \textbf{Backbone / Method} & \textbf{Split Protocol} & \textbf{Key Strengths \& Identified Limitations} \\
\midrule
Traditional CAD Systems (pre-2016) & Hand-crafted ABCD features + SVM & Random / Small private sets & Interpretable features, but fragile under lighting variations, hair artifacts, and multi-class overlap. \\
Esteva et al. (2017) & GoogleNet Inception-v3 transfer learning & Clinician-curated split & Matched dermatologist sensitivity; required $>$129k private images unavailable for open benchmarking. \\
Standard ResNet / DenseNet & ResNet-50 / DenseNet-121 + Cross-Entropy & Naive image-level stratified split & Prone to severe lesion-level data leakage across duplicate HAM10000 views and overconfident minority errors. \\
Vision Transformers (ViT / Swin) & Self-attention patch-based backbones & 5-Fold Cross-Validation & Strong global context modeling, but lacks local translation equivariance and requires heavy regularization on 10k images. \\
\textbf{Proposed JGEC Pipeline} & \textbf{ConvNeXt-Small (IN-22k) + LS + TTA} & \textbf{Strict \texttt{lesion\_id} disjoint} & \textbf{Zero lesion leakage; combines ViT-inspired macro design with CNN inductive bias and calibrated TTA inference.} \\
\bottomrule
\end{tabularx}
\end{table}

\end{document}
